\section{Reproduction entry points and pending Windows work}
\label{app:reproduction}
The report is editable LaTeX. Research/native helpers require a complete
Texloom checkout at the recorded commits, prepared verifier resources and the
declared dependencies. A report-source ZIP alone is not that checkout. Use
fresh synthetic output directories and preserve every prior failed attempt.

\subsection{Required engineering commands}
The repository requires frontend and default-core checks for source changes:
\begin{verbatim}
npm test
cargo test --locked -p texloom-core
\end{verbatim}
The environment in this run uses a project-local Rust toolchain/cache. That
choice and the locked dependency files must be recorded when reproducing a
build. A successful command is a check of the named code/environment, not an
installed-artifact or native-interaction result.

The opt-in native development example is built without launching the app:
\begin{verbatim}
cargo build --locked -p texloom-desktop --features benchmark `
  --example native_reference_development
\end{verbatim}
These examples use PowerShell backtick continuation syntax. The recorded
run used the resulting Windows debug example. Do not substitute an ordinary
or diagnostic GUI executable for this CLI harness.

\subsection{Preparing and executing a development pilot}
The Python research requirements pin Pillow and pypdf. With the complete
recorded source, prepare an unused output directory:
\begin{verbatim}
python research/tools/native_reference_pilot.py prepare `
  artifacts/NEW_DEVELOPMENT_RUN
\end{verbatim}
This records criteria and source/code fingerprints before any experimental
compilation. Evaluation is then explicitly supplied the native runner and
the independent PDF tools:
\begin{verbatim}
python research/tools/native_reference_pilot.py evaluate `
  artifacts/NEW_DEVELOPMENT_RUN `
  --native-runner target/debug/examples/native_reference_development.exe `
  --pdftotext 'PATH_TO_PDFTOTEXT' --pdftoppm 'PATH_TO_PDFTOPPM'
\end{verbatim}
Replace the illustrative paths and continuation syntax for the local shell.
The exact commands, paths and file hashes used by the recorded runs are in
their receipts and logs. The helper rejects overwritten run/output directories
and implementation changes after preparation. Reproducing the first protocol
requires its recorded source/snapshots; current corrected code is not silently
substituted into its old result.

All experimental LaTeX inputs are explicitly named \code{main.tex} and known
supporting \code{.sty} files. Gold builds disable shell escape and automatic
MiKTeX package installation. Never compile \code{.text} files. No user document
is an experimental fixture, and no ablation is allowed to commit a real project.

\subsection{Rebuilding this report}
From \path{research/report/}, compile only \code{main.tex} with an existing
LaTeX installation. The MiKTeX flags used here are:
\begin{verbatim}
pdflatex --disable-installer -no-shell-escape `
  -interaction=nonstopmode -halt-on-error main.tex
\end{verbatim}
Use a fresh ignored output directory when retaining another version, run
BibTeX against that directory with the report source on its bibliography search
path, and repeat LaTeX until references settle. Check process exit statuses,
render the latest PDF and inspect layout. Compilation alone verifies neither
the truth of claims nor visual quality. Other distributions require their
corresponding no-auto-install configuration; do not copy MiKTeX-specific flags
as a universal engine contract.

\subsection{Windows gates requiring version-specific evidence}
\begin{enumerate}
\item Reopen the exact replacement ordinary r5 application with its matching resources,
using disposable synthetic projects; repeat the supported source/signature/
control/Save/Undo/reopen workflows. The target build directory can change between
ordinary and diagnostic builds; select by receipt, not directory name alone.
\item Retain the owner's reported Escape/Alt+Tab cancellation result on the
earlier ordinary copy, and distinguish it from replacement-build lost-capture
and diagnostic observations. Simulated
event tests are separate; the present control API could not inject every
interleaving of held pointer and keyboard input.
\item Test actual DPI changes and multiple-monitor transitions. A zoom/scroll
anchor test is not a Windows display-scale test.
\item Validate clean-account prerequisites, WebView2/setup guidance, missing
compiler/package/font recovery, the current installer and permitted uninstall
procedure. The previously blocked uninstaller launch must not be bypassed.
\item Validate diagnostic boundaries and paired external overhead, then run
the declared timing and single-instance resource workloads. Collect an
unassisted usability task separately.
\end{enumerate}
These steps are pending, not instructions to perform them while the owner is
away. They require the owner's deferred application access/approvals. Separate
independent labels, oracle review and final manifests are also required before
the final evaluation; signing/public distribution remains a later owner decision.
